\documentclass[../cours_1re.tex]{subfiles}
\begin{document}
\setcounter{chapter}{22}
\chapter{Échantillonnage et simulation}
\label{ch:23}
\begin{objectif}
Simuler une variable, observer la fluctuation de moyennes et justifier l'échelle en $1/\sqrt n$.
Prérequis : variables aléatoires et variance. Les inégalités probabilistes sont des compléments.
\end{objectif}
\section{Échantillon et moyenne}
Un échantillon de taille $n$ est ici une famille $X_1,\ldots,X_n$ de variables
indépendantes, de même loi qu'une variable $X$ d'espérance $\mu$ et de
variance $\sigma^2$. Sa moyenne est $\overline X_n=(X_1+\cdots+X_n)/n$.
Une réalisation est une liste de nombres ; l'échantillon théorique est une
famille de variables aléatoires avant l'observation.
\begin{propriete}
$\mathbb E(\overline X_n)=\mu$ et
$\operatorname{Var}(\overline X_n)=\sigma^2/n$.
L'écart type de la moyenne est $\sigma/\sqrt n$.
\end{propriete}
\begin{demonstration}
La linéarité donne $\mathbb E(\sum_iX_i)/n=n\mu/n=\mu$.
L'indépendance donne $\operatorname{Var}(\sum_iX_i)=n\sigma^2$.
Diviser une variable par $n$ divise sa variance par $n^2$, d'où le résultat.
\end{demonstration}
Multiplier $n$ par $4$ divise l'écart type par $2$. Multiplier $n$ par $100$
le divise seulement par $10$. Une dépendance entre observations peut invalider
ce calcul, car les covariances ne sont plus forcément nulles.

\section{Une garantie démontrable}
\approfondissement
Si $Y\ge0$ et $a>0$, on a
$\mathbb E(Y)\ge aP(Y\ge a)$ : dans la somme définissant l'espérance,
les issues où $Y\ge a$ contribuent au moins $a$ fois leur probabilité
et les autres termes sont positifs. C'est l'inégalité de Markov.
L'appliquer à $Y=(X-\mu)^2$ donne, pour $t>0$,
\[
P(|X-\mu|\ge t)\le\frac{\sigma^2}{t^2}.
\]
C'est l'inégalité de Bienaymé--Tchebychev. Pour la moyenne :
\[
P(|\overline X_n-\mu|\ge\varepsilon)
\le \frac{\sigma^2}{n\varepsilon^2}.
\]
Pour tout $\varepsilon>0$ fixé, cette borne tend vers $0$.
On vient de démontrer une loi faible des grands nombres dans ce cadre.
Ce résultat porte sur une probabilité d'écart, pas sur la convergence
monotone d'une liste de moyennes observées.

\section{L'expérience demandée au programme}
Fixer une loi, sa moyenne $\mu$, son écart type $\sigma$, une taille $n$
et un nombre $N$ de répétitions. Simuler $N$ échantillons indépendants,
calculer leurs moyennes, puis la proportion vérifiant
\[
|\overline X_n-\mu|\le\frac{2\sigma}{\sqrt n}.
\]
Pour $\sigma>0$, l'inégalité précédente garantit une probabilité au moins
$3/4$ pour cet événement. Si $\sigma=0$, il se produit avec probabilité $1$.
La proportion simulée fluctue elle-même autour de cette probabilité.

\begin{attention}
La valeur d'environ $95\%$ est liée à une approximation normale appropriée,
pas à une propriété de toutes les lois. Pour $n=1$ et une Bernoulli de paramètre
$0{,}1$, l'intervalle $\mu\pm2\sigma=[-0{,}5;0{,}7]$ contient seulement
la valeur $0$ ; la probabilité est donc $0{,}9$.
Le théorème central limite sera étudié plus tard, avec ses hypothèses.
\end{attention}
Le module Python joint reproduit cette expérience, avec une graine indiquée.
Une graine rend un calcul reproductible ; elle ne prouve pas que le
générateur fournit une indépendance mathématique parfaite.
\begin{exemple}[une expérience effectivement exécutée]
Pour une Bernoulli de paramètre $0{,}3$, $n=100$, $N=2000$ et la graine
$2026$ du module joint, la proportion observée est $0{,}954$.
Le seuil est $2\sqrt{0{,}21}/10\simeq0{,}09165$.
Le résultat décrit cette expérience ; il n'établit pas une loi universelle
à $95\%$. Modifier $n$, le paramètre ou la graine permet d'étudier la fluctuation.
\end{exemple}

\section{Méthode de Monte-Carlo : compter pour mesurer}
Dans le carré unité, on tire un point $(U,V)$ uniformément. La probabilité
de tomber sous la courbe $v=u^2$ est son aire, qui vaut $1/3$
(preuve par intégration au chapitre 25).
La fréquence des réussites estime donc cette aire.
Pour le quart de disque $u^2+v^2\le1$, l'aire est $\pi/4$ ;
quatre fois la fréquence estime $\pi$.

Le tirage uniforme continu utilise un modèle au-delà de l'univers fini
du programme : la probabilité d'une région est définie par son aire.
L'indicatrice de réussite, elle, suit une Bernoulli. Les calculs sur les
indicatrices et leurs moyennes restent donc ceux démontrés plus haut.
Un ordinateur tire sur une grille finie de nombres représentables ;
l'approximation du modèle continu est une hypothèse de simulation.

Pour une aire $p$, l'écart type de sa fréquence sur $n$ points est
$\sqrt{p(1-p)/n}\le1/(2\sqrt n)$, car
$p(1-p)=1/4-(p-1/2)^2\le1/4$.
Pour l'estimateur de $\pi$, il est au plus $2/\sqrt n$.
Ces bornes ne sont pas des erreurs maximales certaines.

\begin{finchapitre}
\section{Lire un résultat numérique}
Une étude reproductible doit donner la loi, les paramètres, la taille,
le nombre de répétitions et la graine. Comparer une erreur observée à une
dispersion théorique, puis refaire plusieurs répétitions.
Un seul résultat proche de la valeur exacte ne mesure pas la précision d'une méthode.
Un histogramme de moyennes peut révéler leur concentration, mais l'aire des
rectangles, et non leur hauteur seule lorsque les classes sont inégales,
représente les fréquences.


\section{Utilité et prolongements}
L'échantillonnage explique pourquoi répéter des mesures réduit certaines
fluctuations et pourquoi le gain de précision est lent.

Les statistiques de licence introduiront estimation et intervalles de confiance.
En master, Monte-Carlo permettra d'approcher des intégrales de grande dimension,
avec des méthodes de réduction de variance ; les hypothèses de dépendance
et de modèle resteront décisives.
\begin{retenir}
Sous indépendance et même loi, l'erreur quadratique moyenne décroît
comme $1/n$, l'écart type comme $1/\sqrt n$. Un résultat simulé et un théorème ont des statuts différents.
\end{retenir}
\end{finchapitre}
\end{document}
